
%%%%%%%%%%%%%%%%%%%%%%%%%%%%%%%%%%%%%%%%%%%%%%%%%%%%%%%%%%%%%
\subsubsection{具体分析}

问题四要求基于建模经验，说明商超还需采集哪些数据并论证其帮助。本问以前三问建模过程中实际暴露的数据局限性为锚点，每条建议均对应一个在Q1--Q3中被具体量化的模型缺陷，通过追溯该缺陷如何影响利润、精度或可行解的丧失，论证采集该项数据对模型改进的必然性。最终按缺陷严重程度与数据可得性排序。

在前三问的建模过程中，识别出以下关键数据局限性：无法区分零销量源于无需求还是缺货（影响需求估计的准确性），折扣回收率$k$的估算精度受折扣原因混杂影响（$k$是模型最敏感参数），Prophet残差仍存在未被模型解释的结构性变动（暗示遗漏了天气、促销等外部变量），部分单品因缺乏上下架时间信息而难以判断其季节性模式，以及价格弹性估计受遗漏变量影响可能存在偏误。

\subsubsection{模型准备}

本问不涉及数学模型，而是以前三问的实际建模经验作为证据来源，与文献中的建议方向交叉验证。相关文献中，文献\cite{Wu2024}在数据处理中揭示了缺货和库存数据的必要性，文献\cite{Cao2024}建议采集天气、节假日和产地数据，文献\cite{Chen2024}建议收集天气、竞争和消费者反馈数据。本文在此基础上，以Q1--Q3的实际数据缺口验证每条建议的边际价值，并按采集成本和预期收益排序。
